\section{Related Work}\label{sec:related-work}

\cnote{C34}{Compress? This is already the shortest section (about 0.35 page) and both ADBIS reviewers asked to \emph{strengthen} it (D6, R4). Cut the filler marked below, but spend the saved lines on the missing references: classical path algebras (Carr\'e, Lehmann, Gondran--Minoux, Mohri), provenance semirings, RPQ engines with GQL path modes, and the group's own CFPQ-in-linear-algebra work (review-ai.md, M6). Net length: unchanged or slightly longer.}
The algebraic approach to path problems has a long history. Kleene's algorithm computes reachability using matrix operations over the Boolean semiring.
This approach has been generalized to arbitrary semirings for solving various algebraic path problems~\cite{graph-using-la}, including path enumeration.
Different query semantics and respective algorithms for regular path queries, including all shortest paths and all simple paths, have been studied in the literature~\cite{martens_et_al:LIPIcs.ICDT.2018.19,Mendelzon1989FindingRS,DBLP:journals/corr/abs-2204-11137}.

\cnote{C35}{Three sentences on GraphBLAS/LAGraph; one is enough (``GraphBLAS~[..] standardizes graph algorithms over arbitrary semirings; LAGraph~[..] collects algorithms built on it''). Also cut ``Recent work has explored \ldots'' below (vague, uncited). About 3 lines.}
The GraphBLAS initiative~\cite{Kepner2016MathematicalFO,brock2021graphblas,10.1145/3577195} established a standardized interface for expressing graph algorithms using linear algebra operations.
The GraphBLAS specification supports arbitrary semirings, enabling the expression of algorithms that go beyond traditional numeric computations.
LAGraph~\cite{lagraph} provides a collection of graph algorithms built on GraphBLAS, including BFS, PageRank, and triangle counting.
Recent work has explored linear algebra approaches for various graph problems.
However, RPQ evaluation in this paradigm remains underexplored, with only a few prior works~\cite{10.1007/978-3-031-43980-3_4,belyanin2150single} addressing this problem.

\cnote{C36}{Each of the three positioning sentences can lose its second clause (about 30\%). About 2 lines.}
Several recent results are closely related to ours but target different settings.
cuRPQ~\cite{park2026curpq} accelerates RPQ and conjunctive RPQ (CRPQ) evaluation on GPUs with a dedicated traversal algorithm and visited-set management; we instead build on general-purpose sparse linear algebra kernels on CPUs, and offloading the same matrix operations to GPUs is a natural next step for our approach.
Abo Khamis et al.~\cite{abokhamis2026acyclic} give output-sensitive algorithms for acyclic CRPQs whose complexity matches that of the corresponding conjunctive queries; this is a theoretical result at the level of CRPQs under reachability semantics, whereas we focus on the practical single-source evaluation of 2-RPQs under path semantics.
LD-RPQB~\cite{ma2025ldrpqb} is a synthetic RPQ benchmark whose data graph follows prescribed path-length distributions; it complements RPQBench~\cite{rpqbench}, which we use, and is a candidate for extending our evaluation.